\documentclass[12pt]{article}
\usepackage[utf8]{inputenc}
\usepackage[T1]{fontenc}
\usepackage{hyperref}
\title{Sukur Language Resources for K-12 Curriculum Development}
\author{Compiled by Claude Sonnet 5.5 (claude-sonnet-5-5) for Andr\'as Kornai}
\date{}
\begin{document}
\maketitle

\section{Introduction}
Sukur (also Sakun, Gemasakun, Sugur; ISO 639-3 syk, Glottocode suku1272) is a Biu--Mandara (Afro-Asiatic) language of Madagali LGA, Adamawa State, Nigeria. Estimates are about 15,000 speakers, up to 27,000 in one source. The Endangered Languages Project rates it Vulnerable, with Hausa gaining ground. No school materials were found.

\section{Primary Educational Resources}
None found. The ELDP collection says it supports a future grammar, dictionary and teaching materials.

\section{Linguistic and Reference Materials}
\begin{itemize}
\item Michael Thomas, \emph{A Grammar of Sakun (Sukur)}, PhD thesis, University of Colorado, 2014. It is the first systematic grammar. \url{https://scholar.colorado.edu/concern/graduate_thesis_or_dissertations/1c18dg03s}
\item Wikipedia article on the language. \url{https://en.wikipedia.org/wiki/Sukur_language}
\item Endangered Languages Project entry. \url{https://endangeredlanguages.com/es/node/97656}
\end{itemize}

\section{Religious and Community Resources}
None found in the search. The SIUE exhibit gives background on the language. \url{https://iris.siue.edu/eng-318/exhibits/show/sakunlanguage/aboutsakun}

\section{Digital and Audio Resources}
ELDP/ELAR collection ``Sakun (Sukur) Language Documentation'' holds an annotated video corpus (about 24 hours reported). \url{https://eldp.access.preservica.com/dk0252/} A DFG project page may relate to the research. \url{https://gepris.dfg.de/gepris/projekt/282128451} A page for the Mainz researcher Friederike Vigeland also appeared. \url{https://www.ifeas.uni-mainz.de/friederike-vigeland-m-a}

\section{Curriculum Development Context}
The language is spoken by the whole community, so intergenerational transmission is not the first concern. Hausa contact is growing. The corpus and grammar make a school curriculum feasible, but the orthography, a lexicon and graded materials would need to be developed with the community. The DFG page and the Vigeland page were not checked for relevance.
\end{document}